\documentclass[10pt, a4paper]{article}

% ==========================================
% PRÉAMBULE
% ==========================================
\usepackage[utf8]{inputenc}
\usepackage[T1]{fontenc}
\usepackage[french]{babel}
\usepackage{geometry}
\usepackage{hyperref}
\usepackage{fancyhdr}
\usepackage{xcolor}
\usepackage{booktabs}
\usepackage{array}
\usepackage{longtable}
\usepackage{enumitem}
\usepackage{titlesec}

\geometry{top=2cm, bottom=2cm, left=2cm, right=2cm}
\setlist{nosep, leftmargin=1.2em}

\hypersetup{
    colorlinks=true,
    linkcolor=black,
    filecolor=magenta,
    urlcolor=blue,
    pdftitle={Liste des competences CARIBOOKS - Alexandre REY},
}

\pagestyle{fancy}
\fancyhf{}
\fancyhead[L]{Liste des Compétences — CARIBOOKS}
\fancyhead[R]{Alexandre REY - ESTIAM}
\fancyfoot[C]{\thepage}

\titlespacing*{\section}{0pt}{10pt}{4pt}

\newcolumntype{L}[1]{>{\raggedright\arraybackslash}p{#1}}

\title{\textbf{Liste des Compétences Travaillées — CARIBOOKS}\\[4pt]
\large\textit{Traçabilité référentiel CDA $\leftrightarrow$ réalisations du projet}}
\author{Alexandre REY \\ E3 DAD — Concepteur Développeur d'Applications (CDA) \\ ESTIAM — Encadrant : Mhand Boufala}
\date{Année scolaire 2025-2026}

\begin{document}
\maketitle
\thispagestyle{fancy}

\noindent{\small Ce document détaille, compétence par compétence, comment chacune a été
mobilisée dans CARIBOOKS. Les preuves associées figurent dans le dépôt de code source et dans
le mémoire de projet (\texttt{main.tex}).}

\section{Bloc 1 — Développer une application sécurisée}

\begin{longtable}{@{}L{6cm}L{9.5cm}@{}}
\toprule
\textbf{Compétence} & \textbf{Mise en œuvre dans CARIBOOKS} \\
\midrule
\endhead
Environnement de développement sécurisé & Secrets via \texttt{.env} (jamais versionnés), garde au démarrage contre une \texttt{SECRET\_KEY} faible ou par défaut (\texttt{test\_secret\_key\_guard.py}) \\
Authentification & Inscription, connexion par mot de passe et via Google Sign-In (OAuth), JWT, rôles utilisateur/administrateur \\
Stockage sécurisé des identifiants & Hashage des mots de passe (bcrypt), champ nullable pour les comptes Google \\
Validation des données entrantes & Schémas Pydantic sur toutes les routes FastAPI \\
Contrôle des permissions & Dépendance FastAPI \texttt{AdminUser} = \texttt{Annotated[..., Depends(require\_admin)]}, isolation inter-utilisateurs testée (\texttt{test\_orders\_cross\_user\_isolation}) ; schémas de saisie distincts selon le demandeur (\texttt{UserSelfUpdate} sans champ \texttt{role} pour \texttt{PUT /users/me}, \texttt{UserUpdate} pour l'administrateur), correctif d'une élévation de privilèges vérifiée puis verrouillée par tests (\texttt{test\_security\_regressions.py}) \\
Prévention des injections SQL & Requêtes paramétrées via SQLAlchemy \\
Prévention des vulnérabilités OWASP & Limitation du taux de requêtes sur l'authentification (\texttt{rate\_limit.py}), protection SSRF sur la récupération d'images distantes (\texttt{test\_image\_service\_ssrf.py}), correction d'une élévation de privilèges (A01 \textit{Broken Access Control}) sur \texttt{PUT /users/me} et suppression des fuites de détails techniques dans les réponses d'erreur (bornes de longueur Pydantic remplaçant des \texttt{DataError} SQL renvoyées en \texttt{500}) \\
Conformité RGPD / nLPD & Export et effacement (anonymisation) des données personnelles \\
Accessibilité de l'interface & Attributs \texttt{alt}, \texttt{aria-label}/\texttt{aria-live}, lien d'évitement, focus clavier visible et tableaux sémantiques sur l'ensemble du site (vitrine et back-office) ; pas d'audit RGAA/WCAG externe ni de test lecteur d'écran réel (limites assumées, cf. mémoire) \\
\bottomrule
\caption{Compétences Bloc 1}
\end{longtable}

\section{Bloc 2 — Concevoir et développer une application multicouche répartie}

\begin{longtable}{@{}L{6cm}L{9.5cm}@{}}
\toprule
\textbf{Compétence} & \textbf{Mise en œuvre dans CARIBOOKS} \\
\midrule
\endhead
Analyse et formalisation du besoin & Cahier des charges, échanges avec la Product Owner \\
Maquettage UI et enchaînement des écrans & Première version des maquettes sous Figma Make (cf. mémoire), écrans finalisés directement en développement au fil des phases \\
Modélisation UML & Cas d'utilisation, séquence, classes, déploiement \\
Modèle de données relationnel & Modélisation Merise (MCD/MLD), base MySQL, scripts de migration versionnés (\texttt{backend/migrations/}) ; sauvegarde/restauration des données de production non automatisée (limite assumée, cf. mémoire) \\
Architecture logicielle en couches & Présentation / Services / Infrastructure (\texttt{backend/presentation}, \texttt{services}, \texttt{infrastructure}) \\
Accès aux données & Couche SQLAlchemy, CRUD générique réutilisable (\texttt{CrudBase}) \\
API REST & Endpoints FastAPI documentés, validés par Pydantic \\
Interface web découplée du back-end & Frontend Next.js / React / TypeScript \\
Intégration de services tiers & OpenLibrary (métadonnées), PostFinance (paiement), Google OAuth (connexion) \\
Cohérence en environnement concurrent & Réservation temporaire de stock (verrou \texttt{with\_for\_update}), libération automatique à expiration ; sérialisation des mutations d'un panier par verrou sur la ligne \texttt{commande} et recalcul du total par lecture verrouillante, corrigeant sous \texttt{REPEATABLE READ} un total erroné et un interblocage InnoDB reproduits en test (\texttt{test\_cart\_integrity.py}) \\
\bottomrule
\caption{Compétences Bloc 2}
\end{longtable}

\section{Bloc 3 — Préparer le déploiement d'une application}

\begin{longtable}{@{}L{6cm}L{9.5cm}@{}}
\toprule
\textbf{Compétence} & \textbf{Mise en œuvre dans CARIBOOKS} \\
\midrule
\endhead
Automatisation du déploiement & Pipeline GitHub Actions : SSH \textit{rolling deploy} vers deux VMs, service \texttt{systemd}, migrations SQL appliquées automatiquement \\
Intégration continue & Pipeline GitHub Actions, tests (pytest + Vitest) à chaque changement \\
Déploiement continu & Backend redéployé automatiquement sur \texttt{VM1} (production) et \texttt{VM2} (secours à chaud) à chaque \textit{push} sur \texttt{main} ; frontend déployé sur Vercel \\
Choix et configuration de l'hébergement & Frontend Vercel, backend/MySQL sur VMs Oracle Cloud, reverse proxy Nginx (TLS) \\
Sécurisation de la production & HTTPS, variables d'environnement séparées du code, base MySQL isolée sur réseau privé (VCN) \\
Plan et campagne de tests & 151 tests automatisés (112 backend + 39 frontend), couverture pytest par module, jeu d'essai détaillé sur les deux parcours les plus représentatifs (commande/paiement, scan ISBN) \\
Vérification de la qualité du code & Ruff (backend) et ESLint (frontend) exécutés en CI à chaque \textit{push}, dépôt intégralement conforme \\
Documentation du déploiement & Procédure documentée et reproductible \\
\bottomrule
\caption{Compétences Bloc 3}
\end{longtable}

\section{Compétences transversales}

\begin{longtable}{@{}L{6cm}L{9.5cm}@{}}
\toprule
\textbf{Compétence} & \textbf{Mise en œuvre dans CARIBOOKS} \\
\midrule
\endhead
Gestion de projet Agile & Flux continu Kanban en solo (backlog Microsoft Planner, cinq grandes phases sur environ 3 mois) \\
Communication avec un commanditaire non technique & Démonstrations à chaque étape marquante \\
Rédaction technique & Mémoire, cahier des charges, documentation de déploiement \\
Gestion de version (Git/GitHub) & Commits structurés, branches de correctifs isolées \\
Esprit critique face aux outils IA & Utilisation supervisée, relecture systématique \\
\bottomrule
\caption{Compétences transversales}
\end{longtable}

\end{document}
